\section{GenConfig}
\label{section:genconfig}

\textbf{GenConfig} is a lightweight web application in \texttt{GenConfig/} for creating, validating, and downloading Proteo JSON configuration files. It implements the same schema and validation rules documented in Section~\ref{section:json-config}.

\subsection{Install and launch}

\begin{lstlisting}[style=bashstyle,caption={Start GenConfig.},captionpos=b]
cd GenConfig
pip install -r requirements.txt
python3 app.py
\end{lstlisting}

Open \url{http://127.0.0.1:5000} in a browser. Optional environment variables: \texttt{CONFIG\_GEN\_HOST}, \texttt{CONFIG\_GEN\_PORT}.

\begin{figure}[htbp]
  \centering
  \includegraphics[width=0.95\textwidth]{GenConfig.png}
  \caption{GenConfig editor (Single mode).}
  \label{fig:genconfig-ui}
\end{figure}

\subsection{Editor workflow}

\begin{enumerate}
  \item On start, the editor loads a \textbf{minimal blank} configuration (one phase, one stage, one group).
  \item \textbf{New}: in Single mode loads a template from \texttt{Codes/test.json} when available; in Multi mode loads a blank complex scaffold.
  \item Edit \textbf{General}, \textbf{Phases} (collapsible, with nested stages), and \textbf{Groups} sections.
  \item Set \textbf{Config file name} for downloads (default \texttt{config.json} or \texttt{complex.json} in Multi mode).
  \item \textbf{Validate}: runs structural, stage, and group checks; applies field cleanup (\textit{sanitization}).
  \item \textbf{Download JSON}: saves the current configuration locally. Download is never blocked; validation messages are advisory.
  \item \textbf{Open file\ldots}: load JSON from disk; auto-detects Multi mode when variant strings contain \texttt{:}.
\end{enumerate}

Counts stay synchronized automatically:
\begin{itemize}
  \item \texttt{Total\_Phases} = number of phase cards.
  \item \texttt{Total\_Resizes} = number of groups $-$ 1.
  \item Each phase's \texttt{Total\_Stages} = number of stages in that phase.
\end{itemize}

\subsection{Single vs Multi mode}

Use the toolbar toggle \textbf{Single} | \textbf{Multi}:

\begin{description}
  \item[Single mode] Edit one concrete configuration. Numeric fields, select boxes for enums, and strategy arrays. Output is a ready-to-run JSON file for \texttt{singleRun.sh}.
  \item[Multi mode] Edit a \textbf{complex} configuration with combinatorial variants. Scalar fields accept strings such as \texttt{"10:20"}. Strategy fields accept \texttt{"0,1:0"} syntax. Validate shows how many valid expanded outputs would be produced. Download saves the complex file for CLI expansion.
\end{description}

\subsection{Combinatorial variant syntax}

\begin{table}[H]
\centering
\small
\begin{tabular}{@{}llp{6.5cm}@{}}
\toprule
\textbf{Delimiter} & \textbf{Meaning} & \textbf{Example} \\
\midrule
\texttt{:} & Cartesian alternative & e.g.\ \texttt{"500000:1000000"} yields two SDR element counts \\
\texttt{,} & Values kept together & e.g.\ \texttt{"0,1:0"} yields \texttt{[0,1]} or \texttt{[0]} \\
\bottomrule
\end{tabular}
\caption{Multi-mode variant syntax (INI-compatible convention).}
\end{table}

\textbf{Rules:}
\begin{itemize}
  \item Structural counts (\texttt{Total\_Phases}, \texttt{Total\_Resizes}, \texttt{Total\_Stages}) cannot use variants.
  \item \texttt{ADR} values 0--100 with variant \texttt{SDR} are percentages (same as legacy INI \texttt{correct\_adr}).
  \item Parallel \texttt{Procs} / \texttt{FactorS} pairs must have matching variant counts.
  \item Expansion skips combinations with duplicate consecutive \texttt{Procs} across groups.
\end{itemize}

Example complex file: \texttt{GenConfig/examples/complex.example.json}.

\subsection{Expanding complex JSON}

From the repository root:

\begin{lstlisting}[style=bashstyle,caption={CLI expansion.},captionpos=b]
python3 Exec/PythonCodes/read_multiple_json.py \
  GenConfig/examples/complex.example.json my_run_
\end{lstlisting}

Writes \texttt{Desglosed-<date>/my\_run\_0.json}, \texttt{my\_run\_1.json}, \ldots{} Each output passes structural and stage semantic validation.

Then execute any concrete file inside an existing allocation:

\begin{lstlisting}[style=bashstyle]
bash Exec/singleRunCostum.sh Desglosed-2026-07-13/my_run_0.json 0 1 0
\end{lstlisting}

\subsection{Strategy legend in the UI}

GenConfig shows spawn and redistribution strategy names (indices \textbf{1} and above; \textbf{0} means clear and is omitted from the legend):

\textbf{Spawn:} 1=Pthread, 2=Single, 3=Intercomm, 4=Multiple, 5=Parallel.

\textbf{Redistribution:} 1=Pthread, 2=Wait sources, 3=Wait targets.

Leave a strategy field empty for the default clear value \texttt{[0]}. Strategy combination constraints apply when \texttt{ADR}~$>$~0 (Section~\ref{subsubsection:strategy-rules}).

\subsection{Related files}

\begin{table}[H]
\centering
\small
\begin{tabular}{@{}ll@{}}
\toprule
\textbf{File} & \textbf{Role} \\
\midrule
\texttt{app.py} & Flask server and API endpoints \\
\texttt{schema.py} & Field metadata, enums, defaults \\
\texttt{validator.py} & Structural validation \\
\texttt{stage\_rules.py} & Stage sanitization and semantics \\
\texttt{restrictions.py} & Group cross-field rules \\
\texttt{expand\_multi.py} & Variant parsing and Cartesian expansion \\
\texttt{static/app.js} & Browser UI logic \\
\texttt{examples/complex.example.json} & Sample multi-config \\
\bottomrule
\end{tabular}
\caption{GenConfig source files.}
\end{table}

See \texttt{GenConfig/README.md} for API details (\texttt{/api/validate}, \texttt{/api/multi/preview}, etc.).
